\usepackage{amsmath,amsthm,amssymb}
\usepackage{tikzit}
\usepackage{tikzscale}
\usepackage{keycommand}
\usepackage{pgfplots}
\usepackage{hyperref}
\usepackage{relsize}

\graphicspath{{.}{./figures/}}

\input{quantum.tikzstyles}
\input{dots.tikzdefs}

\theoremstyle{definition}
\newtheorem{theorem}{Theorem}[section]
\newtheorem{corollary}[theorem]{Corollary}
\newtheorem{lemma}[theorem]{Lemma}
\newtheorem{proposition}[theorem]{Proposition}
\newtheorem{conjecture}[theorem]{Conjecture}
\newtheorem{definition/}[theorem]{Definition}
%\newtheorem{fact}[theorem]{Fact}
\newtheorem{example/}[theorem]{Example}
\newtheorem{examples}[theorem]{Examples}
\newtheorem{example*}[theorem]{Example*}
\newtheorem{examples*}[theorem]{Examples*}
\newtheorem{remark}[theorem]{Remark}
\newtheorem{remark*}[theorem]{Remark*}
\newtheorem{question}[theorem]{Question}
\newtheorem{assumption}[theorem]{Assumption}

\newtheorem*{theorem*}{Theorem}
\newtheorem*{corollary*}{Corollary}
\newtheorem*{lemma*}{Lemma}
\newtheorem*{proposition*}{Proposition}

\newenvironment{example}{
  \renewcommand{\qedsymbol}{$\blacktriangle$}%
  \pushQED{\qed}\begin{example/}
}{\popQED\end{example/}}

\newenvironment{definition}{
  \renewcommand{\qedsymbol}{$\blacktriangle$}%
  \pushQED{\qed}\begin{definition/}}
{\popQED\end{definition/}}

\newcommand{\C}{\ensuremath{\mathbb{C}}}
\newcommand{\eqrule}[1]{\ensuremath{\stackrel{#1}{=}}}

% BRAS AND KETS
\newcommand{\bra}[1]{\ensuremath{\left\langle #1 \right|}}
\newcommand{\ket}[1]{\ensuremath{\left|  #1 \right\rangle}}
\newcommand{\roundket}[1]{\ensuremath{\left|  #1 \right)}}
\newcommand{\braket}[2]{\ensuremath{\langle#1|#2\rangle}}
\newcommand{\ketbra}[2]{\ensuremath{\ket{#1}\!\bra{#2}}}

\newcommand{\catstate}[1]{\ket{\text{cat}_{#1}}}


\newcommand\cbox[1]{\vcenter{\hbox{#1}}}
\newcommand{\emptydiag}{\,\tikz{\node[style=empty diagram] (x) {};}\,}
\newcommand{\scalar}[1]{\,\tikz{\node[style=scalar] (x) {$#1$};}\,}

\newcommand{\boxdiag}[1]{\,\tikz{\node[style=box] (x) {$#1$};}\,}

\newkeycommand{\boxmap}[style=box][1]{\,\tikz{\node[style=\commandkey{style}] (x) {$#1$};\draw(-1.25,0)--(x)--(1.25,0);}\,}

\newkeycommand{\pointmap}[style=ket][1]{\,\tikz{\node[style=\commandkey{style}] (x) at (-0.3,0) {$#1$};\node [style=none] (3) at (0.9, 0) {};\node [style=none] (3) at (-0.9, 0) {};\draw (x)--(0.7,0);}\,}
\newkeycommand{\point}[style=ket,estyle=][1]{\,\tikz{\node[style=\commandkey{style}] (x) at (-0.3,0) {$#1$};\draw [\commandkey{estyle}] (x)--(0.7,0);}\,}
\newkeycommand{\longpoint}[style=ket,estyle=][1]{\,\tikz{\node[style=\commandkey{style}] (x) at (-0.3,0) {$#1$};\draw [\commandkey{estyle}] (x)--(1.2,0);}\,}

\newkeycommand{\copointmap}[style=bra][1]{\,\tikz{\node[style=\commandkey{style}] (x) at (0.3,0) {$#1$};\draw (x)--(-0.7,0);}\,}

\newcommand{\graypointmap}[1]{\point[style=gray ket]{#1}}
\newcommand{\graycopointmap}[1]{\copointmap[style=gray bra]{#1}}

\newkeycommand{\twopoint}[style=ket,estyle=][1]{\,\tikz{\node[style=\commandkey{style}] (x) at (-0.3,0) {$#1$};\draw [\commandkey{estyle}] (x)+(0.25,0.2)--(0.7,0.2);\draw [\commandkey{estyle}] (x)+(0.25,-0.2)--(0.7,-0.2);}\,}

\newkeycommand{\twopointketbra}[style1=bra,style2=ket,style3=ket][3]{\,%
\begin{tikzpicture}
  \begin{pgfonlayer}{nodelayer}
    \node [style=none] (0) at (0, -1.75) {};
    \node [style=none] (1) at (-0.75, 1.75) {};
    \node [style=\commandkey{style1}] (2) at (0, -0.75) {$#1$};
    \node [style=\commandkey{style2}] (3) at (-0.75, 0.75) {$#2$};
    \node [style=\commandkey{style3}] (4) at (0.75, 0.75) {$#3$};
    \node [style=none] (5) at (0.75, 1.75) {};
  \end{pgfonlayer}
  \begin{pgfonlayer}{edgelayer}
    \draw (0.center) to (2);
    \draw (1.center) to (3);
    \draw (4) to (5.center);
  \end{pgfonlayer}
\end{tikzpicture}\,}

\newkeycommand{\twotwoketbra}[style1=bra,style2=bra,style3=ket,style4=ket][4]{\,%
\begin{tikzpicture}
\begin{pgfonlayer}{nodelayer}
    \node [style=none] (0) at (-1.75, -0.5) {};
    \node [style=none] (1) at (-1.75, 0.5) {};
    \node [style=\commandkey{style1}] (2) at (-0.75, -0.5) {$#1$};
    \node [style=\commandkey{style2}] (3) at (-0.75, 0.5) {$#2$};
    \node [style=\commandkey{style3}] (4) at (0.75, -0.5) {$#3$};
    \node [style=\commandkey{style4}] (5) at (0.75, 0.5) {$#4$};
    \node [style=none] (6) at (1.75, -0.5) {};
    \node [style=none] (7) at (1.75, 0.5) {};
\end{pgfonlayer}
\begin{pgfonlayer}{edgelayer}
    \draw (0.center) to (2);
    \draw (1.center) to (3);
    \draw (4) to (6.center);
    \draw (5) to (7.center);
\end{pgfonlayer}
\end{tikzpicture}\,}

\newcommand{\pointcopointmap}[2]{\,\tikz{
  \node[style=bra] (x) at (-0.5,0) {$#2$};\draw (x)--(-1.5,0);
  \node[style=ket] (y) at (0.5,0) {$#1$};\draw (1.5,0)--(y);
}\,}

\newcommand{\graypointcopointmap}[2]{\,\tikz{
  \node[style=gray bra] (x) at (-0.5,0) {$#2$};\draw (x)--(-1.5,0);
  \node[style=gray ket] (y) at (0.5,0) {$#1$};\draw (1.5,0)--(y);
}\,}


\newkeycommand{\pointbraket}[style1=ket,style2=bra,estyle=][2]{\,%
\begin{tikzpicture}
  \begin{pgfonlayer}{nodelayer}
    \node [style=\commandkey{style1}] (0) at (-0.625, 0) {$#1$};
    \node [style=\commandkey{style2}] (1) at (0.625, 0) {$#2$};
  \end{pgfonlayer}
  \begin{pgfonlayer}{edgelayer}
    \draw [\commandkey{estyle}] (0) to (1);
  \end{pgfonlayer}
\end{tikzpicture}\,}

\newkeycommand{\phase}[style=Z phase dot][1]{\,\begin{tikzpicture}
    \begin{pgfonlayer}{nodelayer}
        \node [style=none] (0) at (1, 0) {};
        \node [style=\commandkey{style}] (2) at (0, 0) {$#1$}; 
        \node [style=none] (3) at (-1, 0) {};
    \end{pgfonlayer}
    \begin{pgfonlayer}{edgelayer}
        \draw (2) to (0.center);
        \draw (3.center) to (2);
    \end{pgfonlayer}
\end{tikzpicture}\,}

\newcommand{\npoint}[1]{\,
\begin{tikzpicture}
	\begin{pgfonlayer}{nodelayer}
		\node [style=none] (22) at (0, -1.25) {};
		\node [style=none] (23) at (0, 1.25) {};
		\node [style=none] (24) at (1, 1) {};
		\node [style=none] (25) at (1, 0.5) {};
		\node [style=none] (26) at (1, -1) {};
		\node [style=none] (27) at (0.5, 0) {$\vdots$};
		\node [style=none] (28) at (0, 1) {};
		\node [style=none] (29) at (0, 0.5) {};
		\node [style=none] (30) at (0, -1) {};
		\node [style=none] (31) at (-1.5, 0) {};
		\node [style=none] (32) at (-0.5, 0) {$#1$};
	\end{pgfonlayer}
	\begin{pgfonlayer}{edgelayer}
		\draw (23.center) to (22.center);
		\draw (28.center) to (24.center);
		\draw (29.center) to (25.center);
		\draw (30.center) to (26.center);
		\draw (31.center) to (23.center);
		\draw (31.center) to (22.center);
	\end{pgfonlayer}
\end{tikzpicture}
\,}

\newcommand{\idphase}[1]{\phase[style=Z dot]{#1}}
\newcommand{\redidphase}[1]{\phase[style=X dot]{#1}}
\newcommand{\greenphase}[1]{\phase{#1}}
\newcommand{\redphase}[1]{\phase[style=X phase dot]{#1}}

\newcommand{\redbasisspider}{\point[style=X dot]{}}
\newcommand{\redpibasisspider}{\point[style=X phase dot]{\pi}}

\newcommand{\greenbasisspider}{\point[style=Z dot]{}}
\newcommand{\greenpibasisspider}{\point[style=Z phase dot]{\pi}}

\newcommand{\redmeasspider}{\copointmap[style=X dot]{}}
\newcommand{\redpimeasspider}{\copointmap[style=X phase dot]{\pi}}

\newcommand{\greenmeasspider}{\copointmap[style=Z dot]{}}
\newcommand{\greenpimeasspider}{\copointmap[style=Z phase dot]{\pi}}

\newcommand{\magicspider}{\point[style=Z phase dot]{\frac{\pi}{4}}}



\newcommand{\idwire}{\,%
\begin{tikzpicture}
  \begin{pgfonlayer}{nodelayer}
    \node [style=none] (0) at (-1.5, 0) {};
    \node [style=none] (3) at (1.5, 0) {};
  \end{pgfonlayer}
  \begin{pgfonlayer}{edgelayer}
    \draw (0.center) to (3.center);
  \end{pgfonlayer}
\end{tikzpicture}\,}


% =======================
% = PARALLELAGRAM BOXES =
% =======================

\makeatletter
\newcommand{\boxshape}[3]{%
\pgfdeclareshape{#1}{
\inheritsavedanchors[from=rectangle] % this is nearly a rectangle
\inheritanchorborder[from=rectangle]
\inheritanchor[from=rectangle]{center}
\inheritanchor[from=rectangle]{north}
\inheritanchor[from=rectangle]{south}
\inheritanchor[from=rectangle]{west}
\inheritanchor[from=rectangle]{east}
% ... and possibly more
\backgroundpath{% this is new
% store lower right in xa/ya and upper right in xb/yb
\southwest \pgf@xa=\pgf@x \pgf@ya=\pgf@y
\northeast \pgf@xb=\pgf@x \pgf@yb=\pgf@y

\@tempdima=#2
\@tempdimb=#3

\pgfpathmoveto{\pgfpoint{\pgf@xa - 5pt + \@tempdima}{\pgf@ya}}
\pgfpathlineto{\pgfpoint{\pgf@xa - 5pt - \@tempdima}{\pgf@yb}}
\pgfpathlineto{\pgfpoint{\pgf@xb + 5pt + \@tempdimb}{\pgf@yb}}
\pgfpathlineto{\pgfpoint{\pgf@xb + 5pt - \@tempdimb}{\pgf@ya}}
\pgfpathlineto{\pgfpoint{\pgf@xa - 5pt + \@tempdima}{\pgf@ya}}
\pgfpathclose
}
}}

\boxshape{NEbox}{0pt}{5pt}
\boxshape{SEbox}{0pt}{-5pt}
\boxshape{NWbox}{5pt}{0pt}
\boxshape{SWbox}{-5pt}{0pt}
\boxshape{EBox}{-3pt}{3pt}
\boxshape{WBox}{3pt}{-3pt}
\makeatother

\newcommand{\boxpointmap}[2]{\,%
\begin{tikzpicture}
  \begin{pgfonlayer}{nodelayer}
    \node [style=ket] (0) at (-1, 0) {$#1$};
    \node [style=box] (1) at (0.5, 0) {$#2$};
    \node [style=none] (2) at (1.75, 0) {};
  \end{pgfonlayer}
  \begin{pgfonlayer}{edgelayer}
    \draw (0) to (1);
    \draw (1) to (2.center);
  \end{pgfonlayer}
\end{tikzpicture}%
\,}



\makeatletter

\pgfdeclareshape{cornerpoint}{
\inheritsavedanchors[from=rectangle] % this is nearly a rectangle
\inheritanchorborder[from=rectangle]
\inheritanchor[from=rectangle]{center}
\inheritanchor[from=rectangle]{north}
\inheritanchor[from=rectangle]{south}
\inheritanchor[from=rectangle]{west}
\inheritanchor[from=rectangle]{east}
% ... and possibly more
\backgroundpath{% this is new
% store lower right in xa/ya and upper right in xb/yb
\southwest \pgf@xa=\pgf@x \pgf@ya=\pgf@y
\northeast \pgf@xb=\pgf@x \pgf@yb=\pgf@y

\pgfmathsetmacro{\pgf@shorten@left}{\pgfkeysvalueof{/tikz/shorten left}}
\pgfmathsetmacro{\pgf@shorten@right}{\pgfkeysvalueof{/tikz/shorten right}}

\pgfpathmoveto{\pgfpoint{\pgf@xa - 5pt}{0.5 * (\pgf@ya + \pgf@yb)}}
\pgfpathlineto{\pgfpoint{\pgf@xb - 1.5 * \pgf@shorten@left}{\pgf@ya - 8pt + \pgf@shorten@left}}
\pgfpathlineto{\pgfpoint{\pgf@xb}{\pgf@ya - 8pt + \pgf@shorten@left}}
\pgfpathlineto{\pgfpoint{\pgf@xb}{\pgf@yb + 8pt - \pgf@shorten@right}}
\pgfpathlineto{\pgfpoint{\pgf@xb - 1.5 * \pgf@shorten@right}{\pgf@yb + 8pt - \pgf@shorten@right}}
\pgfpathclose
}
}

\pgfdeclareshape{cornercopoint}{
\inheritsavedanchors[from=rectangle] % this is nearly a rectangle
\inheritanchorborder[from=rectangle]
\inheritanchor[from=rectangle]{center}
\inheritanchor[from=rectangle]{north}
\inheritanchor[from=rectangle]{south}
\inheritanchor[from=rectangle]{west}
\inheritanchor[from=rectangle]{east}
% ... and possibly more
\backgroundpath{% this is new
% store lower right in xa/ya and upper right in xb/yb
\southwest \pgf@xa=\pgf@x \pgf@ya=\pgf@y
\northeast \pgf@xb=\pgf@x \pgf@yb=\pgf@y

\pgfmathsetmacro{\pgf@shorten@left}{\pgfkeysvalueof{/tikz/shorten left}}
\pgfmathsetmacro{\pgf@shorten@right}{\pgfkeysvalueof{/tikz/shorten right}}

\pgfpathmoveto{\pgfpoint{\pgf@xb + 5pt}{0.5 * (\pgf@ya + \pgf@yb)}}
\pgfpathlineto{\pgfpoint{\pgf@xa + 1.5 * \pgf@shorten@left}{\pgf@ya - 8pt + \pgf@shorten@left}}
\pgfpathlineto{\pgfpoint{\pgf@xa}{\pgf@ya - 8pt + \pgf@shorten@left}}
\pgfpathlineto{\pgfpoint{\pgf@xa}{\pgf@yb + 8pt - \pgf@shorten@right}}
\pgfpathlineto{\pgfpoint{\pgf@xa + 1.5 * \pgf@shorten@right}{\pgf@yb + 8pt - \pgf@shorten@right}}
\pgfpathclose
}
}

\makeatother

\pgfkeyssetvalue{/tikz/shorten left}{0pt}
\pgfkeyssetvalue{/tikz/shorten right}{0pt}


\newkeycommand{\kpoint}[style=kpoint,estyle=][1]{\,\tikz{\node[style=\commandkey{style}] (x) at (-0.05,0) {$#1$};\draw [\commandkey{estyle}] (x)--(1.05,0);}\,}
\newkeycommand{\kpointbraket}[style1=kpoint,style2=kpoint adjoint,estyle=][2]{\,%
\begin{tikzpicture}
  \begin{pgfonlayer}{nodelayer}
    \node [style=\commandkey{style1}] (0) at (-0.735, 0) {$#1$};
    \node [style=\commandkey{style2}] (1) at (0.735, 0) {$#2$};
  \end{pgfonlayer}
  \begin{pgfonlayer}{edgelayer}
    \draw [\commandkey{estyle}] (0) to (1);
  \end{pgfonlayer}
\end{tikzpicture}\,}

\newcommand{\kpointbraketconj}[2]{\,%
\begin{tikzpicture}
  \begin{pgfonlayer}{nodelayer}
    \node [style=kpoint conjugate] (0) at (-0.735, 0) {$#1$};
    \node [style=kpoint transpose] (1) at (0.735, 0) {$#2$};
  \end{pgfonlayer}
  \begin{pgfonlayer}{edgelayer}
    \draw (0) to (1);
  \end{pgfonlayer}
\end{tikzpicture}\,}

\newcommand{\freediag}[1]{\,%
\begin{tikzpicture}
  \draw plot[smooth cycle] coordinates {(0.8, 0) (0.5, 0.6) (-0.2, 0.2) (-0.5, 0) (0, -0.7)};
  \node at (0,0) {$#1$};
\end{tikzpicture}\,}